\section{Conclusion}
\label{sec:conclusion}

\noindent This document has discussed how the energy network is a vital utility which enables business as usual and daily activities for modern society, and the effect which climate change
and inclement weather events have on these activities by interrupting supply to the consumer. Also detailed is the financial penalty this has on the DSO and the motivation to
develop and implement solutions which can mitigate the impact of these events and how the FLISR technique is one such solution.
\\
The research undertaken builds on previous related work to meet this need, utilising modern software and communication technologies, in addition to leveraging existing DSO data streams,
to provide such a solution in a fault tolerant way by implementing a FLISR application which can operate both remotely and at the edge of the network by applying edge computing techniques.
By automating this process, the time of detection and resolution of fault events is minimal, as demonstrated by the results gathered. This grants the potential reduction of the financial
impact of fault events, and the level of customers affected. The development of a simulation platform which provides an intuitive interface to create and test fault events on real distribution grids, enables a pathway
to interact with DSO personnel to demonstrate, validate and gather feedback to make continual improvements to the FLISR Edge solution.
\\
Such tools are vital in bridging the gap between academic research and industry, and provide a pathway towards continual development, commercialisation and the application of
the technologies and techniques explored to other use cases in the energy space, such as ancillary services for grid stability, energy management systems at the industrial and domestic level or interacting with
the energy market to enable intelligent energy flexibility to meet the demands on increased RES, distributed energy sources and prosumer activities.

% section conclusion (end)
